\subsection{尖点・接線・共点型の中央切断とKleinian量子化}
\label{subsec:cubic-CY-three-alpha-additive}

\begin{proposition}\label{prop:cubic-CY-three-alpha-additive}
三交換子の $\gamma\ne0$ の $\alpha$ 型で、先頭式が尖点三次曲線、
二次曲線と接線の和、または相異なる共点三直線を定めるとする。
固定出典\cite{AS4CYpinned}の表示で $\gamma=1$ とし、
$H=F_0+aQ+a^2V$ を次の表から取る。
共点型には先行する共点での排除条件を適用している。
\[
\begin{array}{c|l|l|l}
 &F_0&Q&V\\ \hline
\text{尖点}&b^3+c^2d&Abc+Bbd+Cd^2&Db+Ec+Fd\\
\text{接線}&b^2d+bc^2&Abd+Bcd+Cd^2&Db+Ec+Fd\\
\text{共点}&bc(b-c)&A(b^2-bc+c^2)+d^2&Eb+Fc
\end{array}
\]
各包絡 $\mathcal B$ の中心三次は $kL+ka^3$ である。
$L$ の全係数は証明および公開検証器に示す。
H3の必要条件として尖点型と接線型では $C\ne0$ である。
この条件の下で、また共点型では全係数について、
各 $R=\mathcal B/(L+\lambda a^3)$ は整域であり、
$\pi R$ はGabriel--Krull rank-one torsion-freeである。
$U=(R[a^{-1}])_0$ は順に $E_6,D_5,D_4$ 型の
Kleinian特異点のフィルター付きPoisson量子化となる。
\end{proposition}
\begin{proof}
$a$ は中央であり、三つの残る関係は
$[c,d]=H_b$、$[d,b]=H_c$、$[b,c]=H_d$ である。
二次の右辺は対称化する。尖点型では $a<c<b<d$、
他の二型では $a<b<c<d$ で順序単項式へ還元する。
後述の重みで全交換補正の重みが厳密に下がるためこれらは張り、
AS4のHilbert値から基底になる。$a=0$ の商は同じ順序単項式の基底を持つ。

$L$ を、$a$ の指数で分けて $L=L_0+aL_1+a^2L_2$ とする。
尖点型では
\begin{align*}
 L_0={}&b^3-3cb^2+c^2d,\\
 L_1={}&-Ac^2/2+Acb-6Ccd+(-5B/2+12C)b^2+(B-6C)bd+Cd^2,\\
 L_2={}&(-AB/2+3AC-D+E)c+(AB-3AC-B^2/2+D-2F)b\\
       &+(AC+3BC+F)d.
\end{align*}
接線型では
\begin{align*}
 L_0={}&-b^3/3+b^2c+b^2d+bc^2,\\
 L_1={}&(-A/2+B/2-2C/3)b^2+(-2B+2C)bc+Abd\\
       &+(-3B/2+2C)c^2+(B-2C)cd+Cd^2,\\
 L_2={}&(AB/2-AC-B^2/6+D-F)b+(-AB/2+B^2/2-BC+E-F)c\\
       &+(AC+2BC-2C^2+F)d.
\end{align*}
共点型では
\begin{align*}
 L_0={}&b^2c-bc^2,\\
 L_1={}&Ab^2-(A+2)bc-2bd+(A+2)c^2+2cd+d^2,\\
 L_2={}&(E-2A)b+Fc+(A+2)d.
\end{align*}
三つの交換子を順序語基底へ還元すれば $L$ は中央となる。
公開検証器はこの多項式恒等式に加え、$a=0$ の商の
次数 $1,2,3$ の中心次元がそれぞれ $0,0,1$ であることを
全順序語の交換子行列で確認する。
$a=0$ での中央三次を $L_0$ で生成するので、
任意の中央三次から $L$ の適切な倍を引くと $a$ の倍となる。
$a$ の正則性で残る因子も中央であり、次数2と1の中心の像が
零なので反復して $a^3$ の倍を得る。

尖点型・接線型で $C=0$ とする。
$d$ 軸は $R=\mathcal B/(L+\lambda a^3)$ の点対象を定め、shiftに固定される。
$d$ の作用を恒等として $a,b,c$ の非対角成分を $u,v,w$ とおく。
中央三次は点の一次イデアルの平方に含まれるので、
拡大に対する条件は二次関係の一次化だけである。
尖点型では $u'=u,\ w'=w-Bu,\ v'=v+w+w'$、
接線型では $u'=u,\ v'=v+Bu,\ w'=w-v-v'-Au$ である。
全初期値は自己拡大を与え、すべてのtails自己拡大をこの形で尽くす。
従って $\operatorname{Ext}^1$ は次元3、shiftの行列式は1で、
補題\ref{lem:cubic-odd-Serre-determinant}に反する。
純粋な $a^3$ 切断は整域でないため、残る切断は全て $L+\lambda a^3$ と書ける。

重み $(w(a),w(b),w(c),w(d))$ をそれぞれ
$(0,4,3,6),(0,2,3,4),(0,2,2,3)$ とする。
三つの補正項は全て対応する交換子の重みより小さい。
順序基底により $\operatorname{gr}_w\mathcal B=k[a,b,c,d]$ である。
三型の最高重み $12,8,6$ の切断記号は順に
\[
 b^3+c^2d+Ca d^2,\qquad b^2d+bc^2+Ca d^2,\qquad bc(b-c)+ad^2.
\]
これらを $f$ と書く。$L+\lambda a^3$ は中央で、$f$ は正則だから、
$\operatorname{gr}_w R=k[a,b,c,d]/(f)$ となる。
$f$ は $a$ について原始な一次式である。
実際、$a$ の係数は $d^2$ の非零倍、定数係数は $d$ で割れない。
従って $f$ は既約であり、$R$ は整域となる。
さらに $R/(a)$ の随伴環は $k[b,c,d]/(F_0)$ であり、
$d$ は $F_0$ を割らないので $R/(a)$ 上で正則である。
$d$ の冪が生成する右イデアルの連続商は
$R/(a)$ のHilbert級数に $(1-z)$ を掛けた級数をshiftしたもので、
いずれもtailsで非零である。
ゆえに $\pi(R/(a))$ は有限長でない。
これと中央 $a$ の冪による降鎖を使えば、
命題\ref{prop:r13-all-cut-domain-rank}の二段階Serre商の議論により
$\pi R$ はGabriel--Krull rank-one torsion-freeとなる。

$a=1$ とした順序基底にも同じ計算が適用される。
まず $\mathcal B/(a-1)$ の随伴環が $k[b,c,d]$ であり、切断の最高記号は $f(1,b,c,d)$。
よって $\operatorname{gr}_w U=k[b,c,d]/(f(1,b,c,d))$ を得る。
三型で差1の交換子記号はこの式のJacobian Poisson括弧である。
尖点型は $d\mapsto d-c^2/(2C)$ で $Cd^2+b^3-c^4/(4C)$、
接線型は $d\mapsto d-b^2/(2C)$ で $Cd^2+bc^2-b^4/(4C)$ となる。
共点型は $d^2$ と相異なる三方向を持つ二元三次式の和である。
これは順にKleinianの $E_6,D_5,D_4$ の形である。
\end{proof}

この三型で用いた全計算は
\texttt{verify\_cubic\_CY\_three\_alpha\_additive\_center.py} にある。
元の斉次次数による可換曲面 $\operatorname{Proj}(\operatorname{gr}_w R)$ の
境界 $a=0$ は、先頭曲線の滑らかな点では曲線方向の偏微分、
唯一の特異点である $d$ 軸では $f_a$ が非零なので、
全て可換曲面の滑らかな点に含まれる。
この計算は可換随伴環での記述である。
H3の許容性には、量子化アフィン環の有限大域次元と
射影境界での非可換な正則性を引き続き判定する。
